\section{Claudeのテキスト・ウォーターマークをめぐる議論}

Claudeのテキスト出力に機械可読なウォーターマークを埋め込む仕組みが、X上で大きな議論になった。投稿では、生成時のランダム性を決定的なシードへ置き換えることで、文章の品質を損なわず、コピー後も保持できる形で検出可能な特徴を持たせる仕組みが説明されている。EU AI Actへの対応という文脈も共有され、技術そのものだけでなく、規制対応をグローバルに適用することへの是非も論点となった。

Grokの観測では、Anthropic側のFAQを踏まえた説明とともに、「商業上のハンディになるのではないか」「将来の権利主張に使われるのではないか」といった懸念や推測も広がった。生成物の識別技術が、品質、プライバシー、規制対応を同時に扱うテーマとして注目された形である。
\dailyxsource{Thariq (@trq212)}{https://x.com/trq212/status/2088721023223132213}
